\chapter{Introducción}

\section{Contexto y motivación}

El mantenimiento predictivo busca reconocer cambios de comportamiento antes de
que una avería obligue a detener un activo. En maquinaria rotativa, las señales
de vibración permiten observar fenómenos como impactos, desalineación,
desbalanceo o degradación de rodamientos. Sin embargo, transformar esas señales
en una decisión útil no consiste únicamente en entrenar un detector. También es
necesario justificar cómo se prepararon los datos, por qué se eligió un modelo,
qué evidencia respalda una alerta y qué limitaciones tiene el resultado.

Los modelos de lenguaje ofrecen una forma flexible de interpretar esa evidencia
y coordinar tareas especializadas. Su uso plantea, a la vez, una dificultad
metodológica: si el agente responde en texto libre o ejecuta directamente las
transformaciones, resulta difícil separar razonamiento, cálculo y error. Este
trabajo adopta una separación explícita. Los agentes proponen decisiones
estructuradas; los ejecutores deterministas procesan las señales; y todos los
pasos relevantes dejan artefactos que pueden revisarse después.

El interés principal no es, por tanto, obtener una única métrica de detección,
sino observar si agentes locales relativamente pequeños son capaces de
formular hipótesis razonables, utilizar herramientas, revisar resultados y
aprender de experiencias anteriores sin perder reproducibilidad.

\section{Planteamiento del problema}

El problema se aborda desde dos escenarios complementarios. CWRU Bearing Data
Center proporciona un benchmark etiquetado y acotado que permite comprobar el
pipeline y mantener una referencia de regresión~\cite{cwru_bearing_data_center}.
El escenario de investigación principal es NASA IMS Bearing, donde una secuencia
de adquisiciones describe la evolución del activo hasta el final del ensayo
~\cite{nasa_pcoe_repository}. En este segundo caso, el objetivo no debe reducirse
a clasificar ventanas aisladas: interesa estudiar la tendencia del indicador de
salud, la persistencia de las alertas, su tasa antes de monitorización y su
posición retrospectiva respecto al final registrado.

Sobre ambos escenarios se incorpora una memoria RAG\@. Un recuerdo puede ayudar a
repetir una decisión útil, advertir sobre un fallo anterior o señalar que una
estrategia no es comparable con el caso actual. El riesgo es que una memoria
irrelevante, incorrecta o autorreferencial condicione decisiones posteriores.
Por ello, recuperar el fragmento más próximo no es suficiente: el ciclo necesita
gobierno, filtrado, citas verificables y evaluación de su efecto.

\section{Preguntas de investigación}

El trabajo se organiza alrededor de las siguientes preguntas:

\begin{enumerate}[label=\textbf{PI\arabic*:}]
  \item ¿Pueden agentes basados en modelos de lenguaje locales formular y
        justificar decisiones técnicas útiles dentro de contratos estrictos,
        delegando el cálculo en ejecutores reproducibles?
  \item ¿Cómo cambia su razonamiento cuando reciben memoria de decisiones
        anteriores, y pueden citar, adaptar o rechazar esa memoria de forma
        trazable?
  \item ¿Qué mecanismos de filtrado y promoción reducen la reutilización de
        recuerdos incompatibles, no validados o perjudiciales?
  \item ¿Cómo se comporta el sistema en un problema \textit{run-to-failure},
        donde la unidad longitudinal es la secuencia de snapshots y las
        métricas operacionales difieren de una clasificación binaria?
  \item Una vez fijado el protocolo experimental, ¿qué diferencias aparecen
        entre Qwen 3.5 y Qwen 3 bajo las mismas entradas, herramientas,
        contratos y condiciones de memoria?
\end{enumerate}

La última pregunta se plantea como comparación posterior. No se deben atribuir
diferencias al modelo mientras cambien simultáneamente el dataset, la política
temporal o el corpus de memoria.

\section{Objetivos}

\subsection{Objetivo general}

Diseñar y evaluar una plataforma híbrida y auditable para monitorización
industrial \textit{run-to-failure}. La propuesta combina un observador
determinista frecuente con supervisión multiagente periódica o dirigida por
eventos, de modo que los agentes puedan diseñar y auditar políticas de
detección antes de un eventual despliegue industrial. El estudio presta
especial atención a sus hipótesis y al efecto de una memoria RAG gobernada
sobre las decisiones.

\subsection{Objetivos específicos}

\begin{enumerate}
  \item Coordinar agentes especializados mediante un supervisor jerárquico y
        representar sus decisiones con esquemas estructurados.
  \item Separar el razonamiento generativo de los ejecutores que transforman
        datos, entrenan modelos y calculan métricas.
  \item Mantener trazabilidad de entradas, configuraciones, decisiones,
        resultados, revisiones y artefactos.
  \item Implementar un ciclo de memoria que distinga captura, validación,
        promoción, recuperación, uso y auditoría.
  \item Evaluar la relevancia de la memoria y la fidelidad de las citas, además
        de su posible efecto sobre configuraciones y resultados.
  \item Definir una metodología temporal para NASA IMS que separe ventanas
        intra-snapshot de observaciones longitudinales por snapshot.
  \item Implementar un replay histórico causal que permita ensayar la
        separación entre monitorización determinista y revisión multiagente
        sin presentar la simulación como adquisición industrial en tiempo real.
  \item Comparar baselines interpretables y modelos de anomalías bajo un mismo
        protocolo temporal.
  \item Preparar una comparación controlada entre Qwen 3.5 y Qwen 3.
\end{enumerate}

\section{Aportaciones del trabajo}

Las aportaciones del trabajo se agrupan en cuatro bloques. En primer lugar, una
arquitectura híbrida que mantiene el cálculo frecuente en ejecutores
deterministas y hace observable la revisión periódica de los agentes sin
permitirles ejecutar código arbitrario. En segundo lugar, un ciclo RAG orientado
a episodios de decisión y no solo a documentos: conserva alternativas,
evidencia, resultado y condiciones de reutilización. En tercer lugar, un
protocolo de evaluación que separa calidad de recuperación, fidelidad de uso,
cambio de comportamiento y resultado operacional. Finalmente, una aplicación
del enfoque al seguimiento \textit{run-to-failure} de rodamientos mediante
replay histórico causal.

\section{Alcance y limitaciones}

El sistema se ejecuta localmente y emplea modelos disponibles mediante Ollama.
Los agentes no modifican directamente los datos ni generan código para
ejecutarlo. La arquitectura contempla distintos datasets y backends de memoria,
pero el estudio no pretende resolver despliegue multiusuario, adquisición
continua de planta, planificación distribuida o seguridad productiva. La vista
de monitorización reproduce una trayectoria histórica con un reloj acelerado;
su semejanza visual con una consola en vivo no acredita latencia, disponibilidad
ni integración con sensores industriales.

La evidencia actual debe interpretarse con cautela. CWRU valida el flujo
etiquetado y los pilotos sintéticos sirven como pruebas de integración. Sobre
NASA IMS Set~2 oficial se dispone ya de un baseline PCA causal que recorre sus
984 snapshots con partición 20/10/70, calibración separada y evaluación
longitudinal agregada por adquisición. La vista de decisión oculta a limpiador,
estructurador y modelador los campos retrospectivos. Esta única trayectoria
permite estudiar el mecanismo y describir la evolución del score, pero no
aporta etiquetas continuas de degradación, un onset físico oficial ni evidencia
de generalización. Aún es necesario congelar el corpus de memoria y ejecutar
ablaciones agénticas con réplicas.

\section{Estructura del documento}

El capítulo 2 revisa mantenimiento predictivo, agentes y memoria RAG\@. El
capítulo 3 presenta la arquitectura y los contratos. El capítulo 4 define la
metodología de evaluación agéntica, de memoria y \textit{run-to-failure}. El
capítulo 5 resume las decisiones de implementación relevantes, sin reproducir
la estructura interna del repositorio. Los capítulos 6 y 7 describen el diseño
experimental y los resultados, separando pilotos sintéticos de evidencia real.
El capítulo 8 recoge conclusiones, limitaciones y trabajo futuro.
